% Copyright 2026  Ed Bueler

\documentclass[10pt,
               svgnames,
               hyperref={colorlinks,citecolor=DeepPink4,linkcolor=FireBrick,urlcolor=Maroon},
               usepdftitle=false]{beamer}

\usetheme{Madrid}
\usecolortheme{beaver}
\setbeamercovered{transparent}
\setbeamerfont{frametitle}{size=\large}
\setbeamercolor*{block title}{bg=red!10}
\setbeamercolor*{block body}{bg=red!5}

\usepackage[english]{babel}
\usepackage[latin1]{inputenc}
\usepackage{times}
\usepackage[T1]{fontenc}
% Or whatever. Note that the encoding and the font should match. If T1
% does not look nice, try deleting the line with the fontenc.

\usepackage{empheq,bm,xspace}
%\usepackage{minted}
%\usepackage{hyperref}
\usepackage{tikz}

% If you wish to uncover everything in a step-wise fashion, uncomment
% the following command: 
%\beamerdefaultoverlayspecification{<+->}

\newcommand{\bb}{\mathbf{b}}
\newcommand{\bc}{\mathbf{c}}
\newcommand{\bbf}{\mathbf{f}}
\newcommand{\bl}{\bm{\ell}}
\newcommand{\br}{\mathbf{r}}
\newcommand{\bs}{\mathbf{s}}
\newcommand{\bx}{\mathbf{x}}
\newcommand{\by}{\mathbf{y}}
\newcommand{\bv}{\mathbf{v}}
\newcommand{\bu}{\mathbf{u}}
\newcommand{\bw}{\mathbf{w}}

\newcommand{\bzero}{\bm{0}}

\newcommand{\CC}{\mathbb{C}}
\newcommand{\RR}{\mathbb{R}}

\newcommand{\ddt}[1]{\ensuremath{\frac{\partial #1}{\partial t}}}
\newcommand{\ddx}[1]{\ensuremath{\frac{\partial #1}{\partial x}}}
\renewcommand{\t}[1]{\texttt{#1}}
\newcommand{\Matlab}{\textsc{Matlab}\xspace}
\newcommand{\eps}{\epsilon}

\newcommand{\twovect}[4]{\ensuremath{{#1}_{#2} =
                            \begin{bmatrix} #3 \\ #4 \end{bmatrix}}}

\newcommand{\rbullet}{{\color{FireBrick} \bullet}}

\newcommand*\circled[1]{\tikz[baseline=(char.base)]{
            \node[shape=circle,draw,inner sep=2pt] (char) {#1};}}


\title{Integration is a functional}

\subtitle{week 1 calculation}

\author{Ed Bueler}

\institute[]{UAF Math 617 Functional Analysis}

\date{Spring 2026}


\begin{document}
\beamertemplatenavigationsymbolsempty

\begin{frame}
  \maketitle
\end{frame}

\begin{frame}{Outline}
  \tableofcontents[hideallsubsections]
\end{frame}

\AtBeginSection[]
{% do not put toc here this time
}

\section{integration}


\begin{frame}{integral of a continuous function}

\begin{definition}[Riemann's integral]
for a mesh $a = x_0 < x_1 < \dots < x_j < \dots < x_n = b$, and evaluation points $x_j^* \in [x_{j-1},x_j]$, define
  $$\int_a^b f(x)\,dx = \lim_{n\to\infty} \sum_{j=1}^n f(x_j^*)\,(x_j - x_{j-1})$$ 
\end{definition}

\begin{itemize}
\item not the only definition of an integral
\item we will get to Lebesgue's, and cover it lightly!
\end{itemize}

\begin{theorem}
if $f:[a,b] \to \RR$ is continuous then $\int_a^b f(x)\,dx \in \RR$ is well-defined
\end{theorem}
\end{frame}


\begin{frame}{integration example}

Example: compute $\int_0^3 x\, dx$ from the definition

\vspace{60mm}
\end{frame}


\begin{frame}{fundamental theorem of calculus}

\begin{theorem}
if $f:[a,b] \to \RR$ is continuous and $F:[a,b]\to \RR$ satisfies $F'(x)=f(x)$ then
    $$\int_a^b f(x)\, dx = F(b) - F(a)$$
\end{theorem}

\begin{itemize}
\item example: redo integral on last slide

\vspace{30mm}
\end{itemize}

\end{frame}


\begin{frame}{more integration examples}

example: $\displaystyle \int_0^1 \frac{x^k}{1+x^2}\,dx$ for $k=0,1,2,3$

\vspace{60mm}
\end{frame}


\AtBeginSection[] {
    \begin{frame}{Outline}
    \setbeamertemplate{section in toc}[sections numbered]
    \tableofcontents[currentsection]%,hideallsubsections]
    \end{frame}
}


\begin{frame}{title}

\begin{itemize}
\item foo
   \begin{itemize}
   \item[$\circ$] bar
   \end{itemize}
\end{itemize}
\end{frame}


\section{numerical integration}

\begin{frame}{title}

\begin{itemize}
\item foo
   \begin{itemize}
   \item[$\circ$] bar
   \end{itemize}
\end{itemize}
\end{frame}


\section{vector spaces of functions, and norms}

\begin{frame}{title}

\begin{itemize}
\item foo
   \begin{itemize}
   \item[$\circ$] bar
   \end{itemize}
\end{itemize}
\end{frame}


\section{theory this week}

\begin{frame}{title}

\begin{itemize}
\item foo
   \begin{itemize}
   \item[$\circ$] bar
   \end{itemize}
\end{itemize}
\end{frame}



%\begin{frame}{references}
%
%\begin{itemize}
%{\small
%\item D.~P.~Kingma \& J.~Ba (2014). \href{https://arxiv.org/abs/1412.6980}{\emph{Adam: A method for stochastic optimization}}, preprint arXiv:1412.6980
%}
%\end{itemize}
%\end{frame}

\end{document}
